\section{Exploiting Web Services using SQL Injection \mbox{Attacks}}

\drquestion{Q1.1}{Authentication Bypass via SQL Injection}{%
  Exploit the SQL injection vulnerability in the above authentication query to log in to Ted's account without knowing his password, and explain your attack based on your exploited SQL.
}

The code snippet given in the assignment brief for this sub-question, taken
from \texttt{unsafe\_home.php}, is:

\begin{lstlisting}[language=PHP]
$input_uname = $_GET['username'];
$input_pwd = $_GET['Password'];
$hashed_pwd = sha1($input_pwd);
...
$sql = "SELECT id, name, eid, salary, birth, ssn, address, email,
nickname, Password
FROM credential
WHERE name= '$input_uname' and Password='$hashed_pwd'";
$result = $conn -> query($sql);
\end{lstlisting}

\dranswer

\paragraph{Vulnerability analysis:} Both values placed into \texttt{\$sql}
originate from the client and are never validated. \texttt{\$input\_uname} is
read unmodified from \texttt{\$\_GET['username']}, and \texttt{\$hashed\_pwd}
is the SHA-1 digest of \texttt{\$\_GET['Password']}. Hashing the password does
not protect this statement: SHA-1 only transforms the \emph{password} field
into a fixed-length hexadecimal string, but the vulnerability lies in the
neighbouring \texttt{name} column, where \texttt{\$input\_uname} is
concatenated into the query completely unmodified. Whatever the password
field contains, an attacker who controls the \texttt{username} field alone can
still reshape the statement. Both variables are joined into the query by
direct PHP string interpolation, with no parameterised query (an
\texttt{mysqli} prepared statement with bound placeholders) and no input
escaping (\texttt{mysqli\_real\_escape\_string()}) applied to either one:

\begin{lstlisting}[language={[MySQL]SQL}]
WHERE name= '$input_uname' and Password='$hashed_pwd'
\end{lstlisting}

The single quotes around each variable mark where the developer
\emph{intends} the string literal to start and end, but they do nothing to
stop the value itself from containing a quote character. Because the query is
assembled as plain text before MySQL ever parses it, a quote inside
\texttt{\$input\_uname} is indistinguishable from the quote the developer
typed around the variable. This is the defining trait of a string-based SQL
injection in an authentication check: user input is interpreted as part of
the query's \emph{syntax} rather than as inert \emph{data}, so whoever
controls the input can rewrite the logic of the query, including removing the
password predicate entirely.

\paragraph{Exploiting the vulnerability:} I submitted \texttt{Ted'\#} in the
Username field and left the Password field blank, as shown in
Figure~\ref{fig:q1-1-login-payload}. The browser transports the \texttt{\#}
character in the request line as the percent-encoded sequence \texttt{\%23};
the web server decodes this back to the literal \texttt{\#} before PHP ever
reads \texttt{\$\_GET['username']}, so the value that reaches \texttt{\$sql}
is the raw character, not the encoded form shown in the address bar.

\drfigure{figures/q1-1-login-payload.PNG}{The injected username entered with the password field left blank.}{fig:q1-1-login-payload}

With the Password field empty, \texttt{\$hashed\_pwd} is \texttt{sha1('')}.
Substituting \texttt{Ted'\#} for \texttt{\$input\_uname} therefore produces
the raw statement, in which \texttt{\$hashed\_pwd} appears as its computed
SHA-1 digest:

\begin{lstlisting}[language={[MySQL]SQL}]
WHERE name = 'Ted'#' and Password='da39a3ee5e6b4b0d3255bfef95601890afd80709'
\end{lstlisting}

MySQL parses this left to right. \texttt{Ted} selects the intended account
name. The apostrophe immediately after it is the one I injected: it closes
the string literal that the developer opened with \texttt{name= '}, so the
database compares \texttt{name} against the literal \texttt{Ted} only, not
\texttt{Ted'\#}. Parsing then reaches the \texttt{\#} character. In MySQL,
\texttt{\#} begins a single-line comment that extends to the end of the
statement, so everything after it -- the stray closing quote left over from
the original template, the literal \texttt{and Password=}, and the SHA-1 hash
of the blank password field -- is discarded before execution instead of being
treated as SQL. The effective condition MySQL actually evaluates is therefore
reduced to:

\begin{lstlisting}[language={[MySQL]SQL}]
WHERE name = 'Ted'
\end{lstlisting}

This selects Ted's row by name alone; no password predicate participates in
the query at all. In \texttt{unsafe\_home.php}, a non-empty \texttt{\$id} in
the returned row is treated as a successful authentication: the script starts
a session, stores Ted's identity in it, and calls \texttt{drawLayout()} to
render his profile. Figure~\ref{fig:q1-1-ted-login} confirms this outcome:
the application displays the Ted Profile page populated with Employee ID
50000 and Salary 110000, the values recorded for Ted in
\texttt{sqllab\_users.sql}, confirming that the returned row genuinely belongs
to Ted rather than to a different or empty account.

\drfigure{figures/q1-1-ted-profile.PNG}{Ted's profile displayed after the injected login request.}{fig:q1-1-ted-login}

This demonstrates that the login succeeded without knowledge of Ted's correct
password: the SHA-1 comparison that should have enforced it was never
evaluated, because the injected comment removed it from the executed query
before MySQL had a chance to check it.

\drquestion{Q1.2}{Unauthorized Salary Modification via SQL Injection}{%
  Using Ted's session from Q1.1, exploit the same UPDATE statement's SQL injection vulnerability to change Boby's salary to SALARY\_1 without knowing either user's password.
}

The code snippet given in the assignment brief for this sub-question, taken
from \texttt{unsafe\_edit\_backend.php}, is:

\begin{lstlisting}[language=PHP]
$hashed_pwd = sha1($input_pwd);
$sql = "UPDATE credential SET
    nickname='$input_nickname',
    email='$input_email',
    address='$input_address',
    Password='$hashed_pwd',
    PhoneNumber='$input_phonenumber'
WHERE ID=$id;";
$conn->query($sql);
\end{lstlisting}

\dranswer

\paragraph{Vulnerability analysis:} \texttt{\$id} is read from
\texttt{\$\_SESSION['id']}, so an attacker cannot set it directly through a
request parameter -- it is always the ID of whoever is currently logged in,
Ted (\texttt{ID=5}) in this case. However, \texttt{NickName}, \texttt{Email},
\texttt{Address}, \texttt{Password}, and \texttt{PhoneNumber} are all read
unmodified from \texttt{\$\_GET} and concatenated into \texttt{\$sql} with no
parameterised query and no escaping, exactly as in \texttt{unsafe\_home.php}.
Because the single quotes around each of these fields only mark the
developer's \emph{intended} string boundary, an attacker who breaks out of
any one of them can append arbitrary SQL, including a brand-new
\texttt{WHERE} clause, and comment away the original one. This means the
attacker does not need to control \texttt{\$id} at all: redirecting the
\texttt{WHERE} clause is enough to make the whole statement apply to any
row, regardless of who is logged in.

The full snippet above is only executed when the Password field is
non-empty. The actual file wraps it in a condition:

\begin{lstlisting}[language=PHP]
if($input_pwd!=''){
    $hashed_pwd = sha1($input_pwd);
    $_SESSION['pwd']=$hashed_pwd;
    $sql = "UPDATE credential SET nickname='$input_nickname',email='$input_email',address='$input_address',Password='$hashed_pwd',PhoneNumber='$input_phonenumber' where ID=$id;";
}else{
    $sql = "UPDATE credential SET nickname='$input_nickname',email='$input_email',address='$input_address',PhoneNumber='$input_phonenumber' where ID=$id;";
}
$conn->query($sql);
\end{lstlisting}

I left the Password field empty, so the executed statement is built by the
\texttt{else} branch, which omits the \texttt{Password} assignment entirely.
This branch is also never checked for success: \texttt{\$conn->query(\$sql)}
is called and the script unconditionally redirects to
\texttt{unsafe\_home.php} whether the query succeeded or failed, so the
application itself gives no on-page indication of the outcome.

\paragraph{Exploiting the vulnerability:} While logged in as Ted from Q1.1, I
opened Edit Profile and entered the following value in the NickName field,
leaving Email, Address, and Phone Number at their pre-filled (blank) values
and Password empty, as shown in Figure~\ref{fig:q1-2-edit-payload}:

\begin{lstlisting}
', Salary=1605 WHERE ID=2 --
\end{lstlisting}

\drfigure[0.485\linewidth]{figures/q1-2-edit-payload.png}{The injected NickName payload entered on Ted's Edit Profile form.}{fig:q1-2-edit-payload}

Substituting this into the \texttt{else} branch's \texttt{\$sql}, with
Ted's session \texttt{\$id} equal to \texttt{5}, produces the raw statement:

\begin{lstlisting}[language={[MySQL]SQL}]
UPDATE credential SET nickname='', Salary=1605 WHERE ID=2 -- ',email='',address='',PhoneNumber='' where ID=5;
\end{lstlisting}

The leading apostrophe in the payload closes the \texttt{nickname='...'}
literal immediately, giving \texttt{nickname=''} -- a harmless no-op, since
Boby's NickName is already empty in the seed data. Everything after that
quote is now parsed as SQL rather than string content: \texttt{, Salary=1605}
adds a second column assignment, and \texttt{WHERE ID=2} supplies a brand new
condition that targets Boby instead of Ted. The trailing \texttt{--}, followed
by a space, then opens a MySQL end-of-line comment that discards everything
after it -- the stray closing quote left over from the template, the
remaining \texttt{email}/\texttt{address}/\texttt{PhoneNumber} assignments,
and the real \texttt{where ID=5} clause that would otherwise have kept the
update on Ted's own row. The effective statement MySQL actually executes is
therefore reduced to:

\begin{lstlisting}[language={[MySQL]SQL}]
UPDATE credential SET nickname='', Salary=1605 WHERE ID=2
\end{lstlisting}

This updates Boby's row (\texttt{ID=2}) alone; Ted's row is never touched,
because the real \texttt{where ID=5} clause was commented out before MySQL
parsed it.

Since \texttt{unsafe\_edit\_backend.php} gives no visible success or failure
message, I verified the result out of band by logging out and logging back
in as Boby using the same technique as Q1.1 (\texttt{Boby'\#} with the
Password field left blank), which loads Boby's profile directly from his row
in \texttt{credential}. Figure~\ref{fig:q1-2-boby-profile-salary} shows the
result:

\drfigure[0.485\linewidth]{figures/q1-2-boby-profile-salary.png}{Boby's profile after the attack, showing Salary updated to 1605.}{fig:q1-2-boby-profile-salary}

Boby's Salary field now reads \texttt{1605}, matching SALARY\_1 for student
s4041605. Employee ID (20000), Birth (4/20), and SSN (10213352) are unchanged
from \texttt{sqllab\_users.sql}, confirming this is genuinely Boby's row, and
NickName remains blank, consistent with the injected \texttt{nickname=''}
assignment being a no-op. This demonstrates that Ted's authenticated session
was used to modify a different user's row without knowing Ted's password
(already bypassed in Q1.1) or Boby's password at any point.

\drquestion{Q1.3}{Unauthorized Profile Modification via SQL Injection}{%
  Using Ted's session from Q1.1, exploit the same UPDATE statement's SQL injection vulnerability to change Boby's nickname, email, and address to the specified values without knowing either user's password.
}

\dranswer

\paragraph{Vulnerability analysis:} This sub-question exploits the same
\texttt{UPDATE} statement identified in Q1.2, again via the \texttt{else}
branch of \texttt{unsafe\_edit\_backend.php} that runs when the Password
field is left empty:

\begin{lstlisting}[language=PHP]
$sql = "UPDATE credential SET nickname='$input_nickname',email='$input_email',address='$input_address',PhoneNumber='$input_phonenumber' where ID=$id;";
\end{lstlisting}

In Q1.2 I broke out of the \emph{first} field (\texttt{NickName}) and
commented away everything after it, so the earlier, genuine-looking
\texttt{email}/\texttt{address}/\texttt{PhoneNumber} values were never
applied to any row -- only the injected \texttt{Salary} assignment mattered.
For this sub-question I need three fields (\texttt{NickName}, \texttt{Email},
\texttt{Address}) to actually take effect on Boby's row, so injecting through
the first field will not work: anything before the injection point is
discarded by the comment, and anything after it never executes on the
intended row. Instead, I inject through the \emph{last} field,
\texttt{PhoneNumber}, immediately before the real \texttt{WHERE} clause. A
single \texttt{UPDATE} statement has exactly one \texttt{WHERE} clause
governing every \texttt{SET} assignment in it, so replacing that clause at
the last possible point makes \emph{all} of the preceding assignments --
including the legitimate-looking NickName, Email, and Address values -- apply
to whichever row the injected clause names, rather than to Ted's own row.

\paragraph{Exploiting the vulnerability:} While logged in as Ted, I opened
Edit Profile and entered the following values, leaving Password empty so the
\texttt{else} branch above is used, as shown in
Figure~\ref{fig:q1-3-edit-payload}:

\begin{itemize}
  \item NickName: \texttt{Tran Dinh Hai}
  \item Email: \texttt{s4041605@rmit.edu.vn}
  \item Address: \texttt{RMIT}
  \item Phone Number: \texttt{' WHERE ID=2\#}
\end{itemize}

\drfigure[0.485\linewidth]{figures/q1-3-edit-payload.png}{The injected values entered on Ted's Edit Profile form, with the payload in Phone Number.}{fig:q1-3-edit-payload}

Substituting these into \texttt{\$sql}, with Ted's session \texttt{\$id}
equal to \texttt{5}, produces the raw statement:

\begin{lstlisting}[language={[MySQL]SQL}]
UPDATE credential SET nickname='Tran Dinh Hai',email='s4041605@rmit.edu.vn',address='RMIT',PhoneNumber='' WHERE ID=2#' where ID=5;
\end{lstlisting}

The apostrophe in the Phone Number payload closes \texttt{PhoneNumber='...'}
immediately, giving \texttt{PhoneNumber=''} -- a no-op, since Boby's Phone
Number is already blank in the seed data. \texttt{WHERE ID=2} then supplies
a brand-new condition for the statement, and \texttt{\#} opens a MySQL
comment that discards everything after it: the stray closing quote left over
from the template and the real \texttt{where ID=5} clause. Because this
substitution replaces the statement's only \texttt{WHERE} clause rather than
being commented out itself, every assignment before it -- \texttt{nickname},
\texttt{email}, and \texttt{address} -- is retained and now applies to
\texttt{ID=2}. The effective statement MySQL actually executes is:

\begin{lstlisting}[language={[MySQL]SQL}]
UPDATE credential SET nickname='Tran Dinh Hai',email='s4041605@rmit.edu.vn',address='RMIT',PhoneNumber='' WHERE ID=2
\end{lstlisting}

I verified the result the same way as Q1.2: logging out and back in as Boby
(\texttt{Boby'\#}, Password blank). Figure~\ref{fig:q1-3-boby-profile-updated}
shows the result:

\drfigure[0.485\linewidth]{figures/q1-3-boby-profile-updated.png}{Boby's profile after the attack, showing the updated NickName, Email, and Address.}{fig:q1-3-boby-profile-updated}

Boby's profile now shows NickName \texttt{Tran Dinh Hai}, Email
\texttt{s4041605@rmit.edu.vn}, and Address \texttt{RMIT}, exactly as
submitted. Phone Number remains blank, consistent with the injected
\texttt{PhoneNumber=''} assignment being a no-op. Salary still reads
\texttt{1605} from Q1.2, and Employee ID (20000), Birth (4/20), and SSN
(10213352) remain unchanged, confirming this is still genuinely Boby's row.
This demonstrates that Ted's authenticated session was again used to modify
Boby's row -- this time three fields at once -- without knowing Ted's
password or Boby's password at any point.

\drquestion{Q1.4}{Identification of the Stacked-Query Injection Vulnerability}{%
  Identify and explain the vulnerable SQL queries in the two given web pages to display the tasks of the user with the most declared tasks.
}

The two figures below are given by the assignment brief, consisting of vulnerabilities: the Set View Preference page (Figure~\ref{fig:brief-figure2}) and the View Tasks page (Figure~\ref{fig:brief-figure3}).

\drfigure[0.6\linewidth]{figures/figure 2.PNG}{Assignment brief, Figure 2: the Set View Preference page (\texttt{unsafe\_view\_order.php}).}{fig:brief-figure2}

\drfigure[0.6\linewidth]{figures/figure 3.PNG}{Assignment brief, Figure 3: the View Tasks page (\texttt{unsafe\_tasks\_view.php}).}{fig:brief-figure3}

\dranswer

This attack relies on two vulnerable statements working together: one that
lets an attacker store arbitrary text unchecked, and one that later reads
that text back and executes it as part of a new query.

\paragraph{Vulnerable query 1: writing the sort preference (\texttt{unsafe\_view\_order.php}).}

In the \texttt{unsafe\_view\_order.php} snippet below, taken from the lab's Docker image at \texttt{image\_www/Code/}, we can see how the sort preference submitted from Figure~\ref{fig:brief-figure2} is written to the database:

\begin{lstlisting}[language=PHP,aboveskip=0pt]
$newfavourite = $_POST["favourite"] . " ";
$newfavourite.= $_POST["order"];

if($prefid!=""){
    $sql2 = "UPDATE preference SET favourite='$newfavourite' where PreferenceID=$prefid;";
}else{
    $sql2 = "INSERT INTO preference(favourite,Owner) VALUES ('$newfavourite', $id);";
}
...
$conn->query($sql2);
\end{lstlisting}

\texttt{\$\_POST["favourite"]} and \texttt{\$\_POST["order"]} are taken
directly from the request with no validation -- there is no check that
\texttt{favourite} names a real column (\texttt{Name}, \texttt{Hours},
\texttt{Amount}, \texttt{Description}, \texttt{Type}) or that \texttt{order}
is \texttt{asc}/\texttt{desc}. The two are concatenated into
\texttt{\$newfavourite} and written straight into the \texttt{preference}
table's \texttt{favourite} column via string interpolation, with no
escaping. This statement is itself injectable (an attacker could break out
of the \texttt{'\$newfavourite'} literal the same way as in Q1.2/Q1.3), but
its more significant role here is that it lets an attacker persist
\emph{any} string, unchecked, for later use.

\paragraph{Vulnerable query 2: reading and ordering tasks (\texttt{unsafe\_tasks\_view.php}).}

In the \texttt{unsafe\_tasks\_view.php} snippet below, also taken from \texttt{image\_www/Code/}, we can see how that stored value is read back and used to build the View Tasks query shown in Figure~\ref{fig:brief-figure3}:

\vspace{-18pt}
\begin{lstlisting}[language=PHP,aboveskip=0pt]
$sql = "select favourite from preference where owner=$id limit 1";
...
$favourite = $json_a[0]['favourite'];
if($favourite==""){
    $favourite = "hours desc";
}

$sql1 = "select tasks.Name as taskname, credential.Name as ownername, tasks.Hours,tasks.Amount,tasks.Description,tasks.Type
from tasks, credential where tasks.owner=credential.ID and tasks.owner=$id "
. "order by tasks." . $favourite ;

$queryList = explode(";",$sql1);

foreach($queryList as $subquery){
    if(strlen($subquery)>0) {
        if (!$tempResult = $conn->query($subquery)) {
            die('There was an error running the query [' . $conn->error . ']\n');
        }
        ...
    }
}
\end{lstlisting}

The stored \texttt{favourite} value is read back from \texttt{preference}
and concatenated, unquoted and unescaped, directly after the literal text
\texttt{order by tasks.} to build \texttt{\$sql1}. Because it is never
checked against a whitelist of column names, an attacker who controls its
content controls part of the SQL text itself, not just a data value. Worse,
the script then calls \texttt{explode(";", \$sql1)} and executes
\emph{every} non-empty resulting fragment as its own independent query via
\texttt{\$conn->query(\$subquery)} inside the loop. This is a stacked-query
mechanism built into the application: if \texttt{favourite} contains a
\texttt{;} followed by a second complete SQL statement, that second
statement is executed exactly as if the attacker had submitted it directly.

\paragraph{Why this is exploitable: the complete data flow:} An
authenticated user (an attacker needs no special privilege beyond a normal
session) submits the Set View Preference form with crafted
\texttt{favourite}/\texttt{order} values instead of the intended dropdown
selection and \texttt{asc}/\texttt{desc} text. Query~1 persists this
crafted string verbatim into that user's row in \texttt{preference}, with no
validation. When the same user opens View Tasks, Query~2 reads that string
back out of the database and splices it, still unvalidated, into a new
\texttt{SELECT} immediately after \texttt{order by tasks.}. Because the
resulting statement is split on \texttt{;} and each fragment is executed
separately, a \texttt{favourite} value of the form
\texttt{<valid column>; <second SELECT statement>} causes the second
statement to run as a genuinely separate query against the same database
connection -- for example, one that replaces the intended restriction
\texttt{tasks.owner=\$id} with \texttt{tasks.owner=userIdMaxTasks()}, so it
returns another user's tasks instead of the attacker's own. Query~1 is the
injection point (no input validation on write); Query~2 is where the stored
payload is later parsed apart and executed (no output validation on read,
plus explicit multi-statement execution). Neither vulnerability alone
completes the attack -- Q1.5 exploits both together.

\drquestion{Q1.5}{Cross-User Data Disclosure via Stacked-Query Injection}{%
  Exploit the vulnerabilities identified in Q1.4 to make View Tasks display the tasks of the user with the most declared tasks.
}

\dranswer

\paragraph{Step 1: submit the payload:} This attack needs both vulnerable
statements from Q1.4 to work together. Query~1 is where I plant the
malicious text; Query~2 is where that text later gets executed. While logged
in as Ted, I opened Set View Preference and filled in the form as shown in
Figure~\ref{fig:q1-5-preference-payload}:

\begin{itemize}
  \item Sort field dropdown (\texttt{favourite}): \texttt{Hours}
  \item Asc or Desc field (\texttt{order}):
\end{itemize}

\begin{lstlisting}
;
select tasks.Name as taskname, credential.Name as ownername, tasks.Hours, tasks.Amount, tasks.Description, tasks.Type
from tasks, credential
where tasks.owner=credential.ID and tasks.owner=userIdMaxTasks()
\end{lstlisting}

\drfigure{figures/q1-5-preference-payload.png}{The Set View Preference form with \texttt{Hours} selected and the injected statement in the Asc or Desc field.}{fig:q1-5-preference-payload}

Picking a real column name (\texttt{Hours}) for the dropdown keeps the first
half of the eventual statement valid SQL. The rest of the attack, a
complete second SQL statement, is typed straight into the free-text
\texttt{order} box. That field has no character restrictions and the form
submits it as an ordinary POST request, so no proxy or browser console was
needed to send it.

\paragraph{Step 2: how the payload is stored:} Query~1 combines the two
submitted fields into one string:

\begin{lstlisting}
$newfavourite = $_POST["favourite"] . " ";
$newfavourite.= $_POST["order"];
\end{lstlisting}

With the values above, this builds:

\begin{lstlisting}
Hours ;
select tasks.Name as taskname, credential.Name as ownername, tasks.Hours, tasks.Amount, tasks.Description, tasks.Type
from tasks, credential
where tasks.owner=credential.ID and tasks.owner=userIdMaxTasks()
\end{lstlisting}

Ted already had a row in \texttt{preference}, so this string is written into
it unchanged by Query~1's \texttt{UPDATE} branch, with no check on its
content.

\paragraph{Step 3: how the payload is executed:} When View Tasks next loads
for Ted, Query~2 reads that same string back out of the database as
\texttt{\$favourite} and appends it to the end of a new query it is
building. With Ted's session \texttt{\$id} equal to \texttt{5}, the
resulting \texttt{\$sql1} is (shown here broken across lines for
readability; PHP and MySQL both treat it as one continuous string, with no
real line breaks):

\begin{lstlisting}[language={[MySQL]SQL}]
select tasks.Name as taskname, credential.Name as ownername, tasks.Hours, tasks.Amount, tasks.Description, tasks.Type
from tasks, credential
where tasks.owner=credential.ID and tasks.owner=5
order by tasks.Hours
;
select tasks.Name as taskname, credential.Name as ownername, tasks.Hours, tasks.Amount, tasks.Description, tasks.Type
from tasks, credential
where tasks.owner=credential.ID and tasks.owner=userIdMaxTasks()
\end{lstlisting}

The semicolon in the middle is the one I typed in the \texttt{order} field.
Query~2 then runs \texttt{explode(";", \$sql1)}, which cuts this single
string into two separate pieces wherever a semicolon appears, and its
\texttt{foreach} loop sends each non-empty piece to the database as its own,
independent query. The table below explains what each of the two resulting
queries does:

{\footnotesize
\begin{tabularx}{\linewidth}{lXX}
\toprule
\textbf{Query} & \textbf{Restriction on \texttt{tasks.owner}} & \textbf{What it returns} \\
\midrule
1: the original, legitimate half & \texttt{tasks.owner=5} -- Ted's own session ID, placed there by the application itself & Ted's own declared tasks, sorted by Hours \\
2: the half I injected & \texttt{tasks.owner=userIdMaxTasks()} -- a function call I supplied, with no connection to who is logged in & The declared tasks of whichever user currently has the most tasks, joined with \texttt{credential} so that user's name is shown too \\
\bottomrule
\end{tabularx}
\par}

\paragraph{Step 4: the result:} Figure~\ref{fig:q1-5-tasks-view-two-tables}
shows the View Tasks page after this request. It renders two separate
tables, one for each query above. The first table lists Ted's own two tasks
(\textit{Back-end dev1}, 30 hours; \textit{Front-end dev2}, 40 hours). The
second table lists four tasks, and every one of them has Task Owner
\texttt{Alice} (\textit{On-site client requirement collection},
\textit{Business Analysis}, \textit{Demo App to client}, \textit{Design
System Architecture1}).

\drfigure{figures/q1-5-tasks-view-two-tables.png}{View Tasks after the attack, showing Ted's own tasks and, in the second table, Alice's tasks returned by the injected query.}{fig:q1-5-tasks-view-two-tables}

The second table is the direct output of Query~2's injected half, so
\texttt{userIdMaxTasks()} resolved to Alice's user ID in the database at the
time of this request. This matches \texttt{sqllab\_users.sql}, where Alice
(\texttt{ID=1}) is declared four tasks in the original seed data, more than
any other user.

\paragraph{Conclusion:} An attacker-controlled string, stored with no
checks by Query~1, was later cut apart and run as a second, independent SQL
statement by Query~2. That second statement ignored the restriction that
should have limited results to the logged-in user and instead returned the
declared tasks of the user with the most tasks overall, without that user's
identity or credentials ever being entered anywhere in the attack. This is
the complete attack the question asks for.
